% AI-written, AI cross-reviewed; no human sign-off; not a claim that RH is proved
% Verbatim sections of local blueprint v13. See README.md for macros, edition and source spans.

% blueprint/ch_P_r12.tex:224-585
\subsection{The two standard Weyl lattices are exact duals}\label{r12:weyl}

Let $K$ be the number field with involution, $R=\mathcal O_K$,
$\iota$ a finite index set, $\lambda_\tau$ dominant weights, and
$x\in VV(K,\lambda)$. For the normalized pairing $P$ at the identity
Gram matrices the two claims are
\[
  P(\Delta,x)\subseteq R\ \Longleftrightarrow\ x\in\nabla,
  \qquad P(\nabla,x)\subseteq R\ \Longleftrightarrow\ x\in\Delta .
\]
Here $\Delta=\mathrm{delStd}$ and $\nabla=\mathrm{nabStd}$.
The condition $x\in VV$ is essential for coordinate testing.
The target is \path{GlobalTheory.T.standard_weyl_lattices_fock_dual}
(\path{StandardWeylLatticesFockDualStatement.lean:14--15}). Its inclusive
cone has 5,730 core lines. References to \path{StatementA.lean},
\path{StatementC.lean}, \path{StatementE.lean}, and the tableau files
in this subsection are in \path{Formalism/GradedWeylDivisibility/}
and its \path{FockSpaceTableaux/} subdirectory. Other bare filenames refer
to \path{Development/AbelianVarietyTheory/Proofs/FoundationProofs/}.

\paragraph{Recovering a polynomial coefficient by testing a Weyl word.}
Start over $\mathbb Q$ with one partition $\nu$, $n\le r$, and the
space $V_\nu$ of polynomials of column multidegree $\nu$ invariant
under adding an earlier column to a later column, with an indeterminate
scalar. Let $h$ be the product of the prescribed powers of leading
principal minors and $N=(h,h)_F>0$. Write
$B=(\cdot,\cdot)_F/N$. The lattice $\nabla_{\mathbb Z}$ consists of
the integral polynomials in $V_\nu$ (the defining coefficient identities
descend along $\mathbb Z\hookrightarrow\mathbb Q$);
$\Delta_{\mathbb Z}$ is the integer span of the divided row words $Ah$.
See \path{dualWeyl}, \path{rowGen}, \path{hw}, \path{applyWord},
\path{weyl} (\path{StatementA.lean:24--73}) and \path{bDual}
(\path{StatementE.lean:21--45}).

Row substitutions preserve column totals and commute with column substitutions. Operators on a polynomial occur in the reverse order to successive row substitutions on its frame. A divided row operation selects the coefficient of $t^k$ in a row
addition. Its adjoint for the factorial monomial pairing interchanges
the two rows. On a monomial this is the factorial identity in
\S\ref{r11:fock}; reversing the order of a product gives $A^t$.
The evaluation identity $B(h,H)=H(I)$ on $V_\nu$ follows by checking
products of initial-column minors. Each such product is row homogeneous.
A product with the row weight of $h$ must, successively, use row $i$ in
every minor of size at least $i$ (one-based rows). Such a product is
$\pm h$ or zero. All other row weights vanish both in the pairing against
$h$ and on evaluation at $I$. The spanning assertion for these products
also follows by the leading-term elimination below, over $\mathbb Q$.
This proves $B(A^t h,F)=(AF)(I)$; for integral $F$ the right side is
integral. See \path{fockQ_hw} (\path{LowerTriOfInterlaced.lean:135--169}),
\path{fockQ_applyWord} (\path{FockQWordRowActAdj.lean:8--18}) and
\path{nablaEqDualWeylOf} (\path{ApplyWordMemDualWeyl.lean:110--180}).

Suppose conversely that $F\in V_\nu$ and $B(\Delta_{\mathbb Z},F)$
is integral. For a lower triangular exponent $e=(e_{as})$ with the
prescribed column totals, set the diagonal variables to one, the
strictly upper variables to zero and leave all strict lower variables
independent. The coefficient at $e$ below the diagonal in the resulting
polynomial is $[X^e]F$: the column totals recover the missing diagonal
powers. This same coefficient is $(A_e F)(I)$. Indeed start at $I$
and install below-diagonal variables column by column from right to left.
For position $(a,s)$ add an indeterminate multiple of row $s$ to row
$a$ and extract its power $e_{as}$ in the substitution on $F$.
At this stage row $s$ is still the $s$th coordinate row: all previously
installed entries have column at least $s$ and row strictly larger
than their column. The substitution therefore installs precisely
the variable at $(a,s)$; repeated coefficient extractions recover the
coefficient claimed. This gives the source's ordered word under the
fold-right convention. Exponents of the wrong column totals have zero
coefficient. Thus every lower triangular coefficient is an integer,
by applying the preceding identity to $A_e^th$.
See \path{coeff_single}, \path{map_rowGen}, \path{eval_applyWord},
\path{wordOf_eq}, \path{coeffWord}
(\path{UnitriMonomialOfLowerTri.lean:36--58,60--88,93--131},
\path{UnitriMonomialOfLowerTri.lean:141--156,221--236}).

\paragraph{Leading exponents and integral straightening.}
Take the lexicographic order comparing first by column, then by row.
The leading exponent $m$ of a nonzero $F\in V_\nu$ satisfies, at
each pair of consecutive columns $l,j$ and each row $a$,
\[
       \sum_{b\le a}m_{bj}\ \le\ \sum_{b<a}m_{bl}.                 \tag{*}
\]
Here is the two-column calculation explaining the strict endpoint.
Keep the component with the same exponents outside those columns as
$m$ and with the same combined degree in the two entries of each row.
This retains the leading term, and substitution
$X_{bj}\mapsto X_{bj}+tX_{bl}$ preserves this grading. The transformed
component has $t$-degree zero since the original polynomial is invariant.

More generally, give this transformed polynomial a degree allowance
$s$. At the first remaining row write its variables as $X,Y$ and its
degree there as $d$. Expand $\sum_{k=0}^K X^kY^{d-k}w_k$,
where $w_K\ne0$ and the $w_k$ are independent of $X,Y$.
The column order selects $k=K$ for the leading term. Put $e=d-K$.
After substitution split off $(Y+tX)^e$; the residual factor is
\[
       L=\sum_{k=0}^{K} X^k(Y+tX)^{K-k}\,w_k(t),
\]
where $w_k(t)$ denotes substitution in all other rows.
Substitution is injective (put $t=0$); the polynomial ring is a
domain, so $e\le s$ and $\deg_t L\le s-e$. Putting $X=1$, $Y=-t$
recovers $w_K(t)$, whose $t$-degree is at most $s-e+K$ since $L$
has total $X,Y$-degree $K$. Now remove that row from $w_K$.
At the first row $e\le s$; at every subsequent cut, the remaining
right-column total is at most $s-e+K$ plus the preceding left-column
total in the remaining rows, by induction. Restoring the first row
adds $e$ on the left of the desired inequality and $K$ on the right.
This proves (*) with allowance $s$, with the empty row set as base
case. Apply it to the initial allowance zero. Sources:
\path{InterlacedAt}, \path{pairW} (\path{StatementC.lean:34--46}),
\path{leadInterlaced_of} (\path{LeadInterlaced.lean:30--62}),
\path{lctwo_main} (\path{TwoColumn.lean:47--208}) and
\path{peelDegBound_of} (\path{PeelDegBound.lean:119--181}).

In column $s$ make a sorted list $L_s$ containing row $a$ with
multiplicity $m_{as}$. Its length is $\nu_s$. By (*) the $k$th entry
of $L_s$, whenever present, is strictly larger than the $k$th entry
of $L_{s-1}$: apply (*) with $a$ that $k$th entry. For each $k$ take
the initial-column minor with rows the $k$th entries of these lists,
in column order. Their columns form an initial segment since $\nu$ is
a partition. The row indices strictly increase. The diagonal
permutation gives the unique leading monomial of this minor, with
coefficient one: the first column where another permutation differs
uses a larger row index. The product $D_m$ of the minors is therefore
monic with leading exponent $m$, column totals $\nu$, and integer
coefficients. The minors are invariant under the permitted column
additions. This constructs the required tableau, including the empty
product for the zero exponent. In particular $m$ is lower triangular,
since the $s$th row in any strictly increasing list of row indices
is at least $s$.
See \path{sSYTExists_of} (\path{SSYTExists.lean:11--62}),
\path{colLexInitMinor_of} (\path{ColLexInitMinor.lean:22--52}) and
\path{columnLexBideterminant_of_initialMinor}
(\path{ColumnLexBideterminant.lean:11--27}).

Subtract $([X^m]F)D_m$. The leading exponent falls and every lower
triangular coefficient stays integral. There are finitely many
exponents at the fixed column totals. Iteration ends at zero and
expresses $F$ as an integer sum of the $D_m$. With arbitrary coefficients
the same elimination over a field proves the spanning assertion used
above. Hence $\Delta_{\mathbb Z}^{\#}=\nabla_{\mathbb Z}$ in
$V_\nu$, where $\#$ denotes the dual for $B$.
This proves the converse in \path{nablaEqDualWeylOf}; see also
\path{GlobalTheory.WeylInt.lowerTriIntegral}
(\path{LowerTriIntegral.lean:41--82}).

\paragraph{The other dual and the number-field coefficients.}
The lattice $\Delta_{\mathbb Z}$ is contained in the finite free
coefficient module at the fixed multidegree. It therefore has a finite
integer basis by \path{Submodule.basisOfPid}
(\path{Mathlib/LinearAlgebra/FreeModule/PID.lean}:300--302).
Its rational span is all of $V_\nu$: a vector
orthogonal to it has every lower triangular coefficient zero by word
extraction, whereas a nonzero vector has a lower triangular leading
exponent. The restricted Fock form is positive definite over $\mathbb Q$.
Let $d_i$ be an integer basis of $\Delta_{\mathbb Z}$. Its
$B$-dual basis $b_i$ exists over $\mathbb Q$ by inverting the Gram
matrix. Testing a rational vector against the $d_i$ reads off precisely
its coordinates on the $b_i$. Thus $\nabla_{\mathbb Z}$ has basis
$b_i$. Testing on the $b_i$ gives back $\bigoplus_i\mathbb Z d_i$.
This proves both dualities over $\mathbb Z$ without any unimodularity
assumption on an individual Gram matrix. See \path{dual_dual_of_isLattice},
\path{normB_nondeg}, \path{weylEqDualNablaOf}
(\path{ApplyWordMemDualWeyl.lean:182--217},
\path{ApplyWordMemDualWeyl.lean:298--324}).

Apply this construction at each $\tau$ to $\nu=\lambda_\tau.\mathrm{natW}$.
The divided words commute with coefficient extension, so the $d_i$
span the Weyl lattice over $R$. The leading-term argument over the
domain $R$ expresses each integral dual Weyl polynomial as an $R$-sum
of the same integral bideterminants. They belong to the integral span
of the $b_i$ already over $\mathbb Z$, and the $b_i$ retain their
defining invariance over $R$. Thus these two rational bases also generate
the two lattices defined over $R$. Tensor over the finite set of weights.
On products of polynomials in the disjoint blocks, the factorial monomial
form factors: the monomial factorials and the coefficients both multiply.
The denominator of $P$ is the product of the single-block highest-weight
norms, so the two product bases pair by $\delta_{IJ}$, in either order.
They span the prescribed coefficient space over $K$ and their integral
spans are the two standard lattices. Indeed the product polynomials
give the pure tensors generating those spaces, and duality of the
displayed families gives independence. The involution fixes the rational
basis coefficients and preserves $R$ because it preserves integrality
over $\mathbb Z$. Testing $x$ on either first-argument basis recovers
its second-argument coordinates on the opposite basis; integral sums
of basis elements pair integrally. This proves both biconditionals.
See \path{dualWeyl_toK}, \path{weyl_toK}, \path{exists_dual_bases}, and
\path{std_dual} (\path{StandardWeylLatticesFockDual.lean:129--139},
\path{StandardWeylLatticesFockDual.lean:145--235},
\path{StandardWeylLatticesFockDual.lean:237--307}).

\paragraph{Localization of the transported lattices.}
The preceding argument concerns the standard models. At a finite
prime $w$ write $R_w$ for the localization of $R$ inside $K$.
Localize either standard lattice $T$. Its defining integrality test is
independent of the local frame: two bases of the same full lattice
differ by a matrix in $\mathrm{GL}_n(R_w)$.
For $\nabla$ an integral basis substitution with integral inverse
preserves integral polynomials, column weights and right column
invariance. For $\Delta$, an elementary row addition acts by the sum
of its divided row operations multiplied by the corresponding powers
of its integral parameter, a finite sum on each polynomial. Each
divided operator prepends a letter to a Weyl generator. Diagonal units
act on a word by a unit since $h$ is row homogeneous and every divided
operation shifts the row weight by its prescribed amount. These
matrices generate $\mathrm{GL}_n(R_w)$ by elimination with unit
pivots and row interchanges (the latter are products of elementary
matrices and diagonal units). The determinant twist is a unit even
at a negative exponent. This gives frame independence for both
lattices, including products of blocks.
See \path{isBasisAt_inv_mul}, \path{indep_of_stable},
\path{nabIndep_and_delIndep}
(\path{NabIndepAndDelIndepLocalConditions.lean}:21--43) and the two
stability inputs \path{GlobalTheory.T.polyAction_mem_locP_nabStd_of_isBasisAt}
and \path{GlobalTheory.T.polyAction_mem_locP_delStd_of_isBasisAt}
(\path{ActPMemLocPNabStdOfIsBasisAtStatement.lean}:14--15;
\path{ActPMemLocPDelStdOfIsBasisAtStatement.lean}:14--15).

Given a full lattice $M_\tau$ in each block, choose a local basis
matrix $B_{w,\tau}$ over $K$ at $w$: $R_w$ is a DVR and the finitely
generated torsion-free localization is free of full rank
(\path{GlobalTheory.LatBasisAt.basisAtEx},
\path{LatticeHasBasisAtEveryFinitePlace.lean}:66--105). Write
$T_w(B_w)=\mathrm{act}(B_w)(T\otimes_R R_w)$ and
$P_T=\bigcap_w T_w(B_w)$ in the rational coefficient space.
The equality $(P_T)_w=T_w(B_w)$ merits a denominator argument.
The inclusion from the left follows by localizing the defining test.
For $x\in T_w(B_w)$, away from finitely many primes all $M_\tau$
have identity frame (clear denominators in both inclusions between
two full finitely generated lattices); also $x\in T_v(1)$ away from
finitely many primes since $T$ spans the coefficient space. At each
exceptional $v\ne w$, choose $k_v$ so $v^{k_v}x\subseteq T_v(B_v)$
by clearing local denominators. Choose $s_v\in v^{k_v}\setminus w$,
which exists because distinct maximal ideals are comaximal.
The product of the finitely many $s_v$ belongs to $R\setminus w$
and sends $x$ into every $T_v(B_v)$. Thus $x\in(P_T)_w$. With a
single global frame the same intersection gives precisely the transported
global lattice, since if $x\notin T$, the proper ideal
$\{a\in R:ax\in T\}$ lies in some maximal ideal; there $x\notin T_w$.
Finally, a single local test places $x$ in the $K$-span of the
transported lattice and hence in $VV$, which is stable under changing
frame (row substitution commutes with the defining right-column identities).
These prove the three assertions in
\path{GlobalTheory.T.transportedLattice_localization_free_le_coefficientSpace}
(\path{TransportedLatticeLocalizationFreeLeCoefficientSStatement.lean:14--15});
see \path{exists_isBasisAt_one}, \path{transLat_colLat}, \path{locP_transLat}
and \path{nabDelLoc_free_leVV}
(\path{TransportedLatticeLocalizationFreeLeCoefficientS.lean:91--125},
\path{TransportedLatticeLocalizationFreeLeCoefficientS.lean:139--234}).

\paragraph{Hermitian duality after transport.}\phantomsection\label{r13:weyl-dual}
Let $R=\mathcal O_K$ for a number field $K$ with involution. Fix a
finite set $\iota$, a rank $n$, dominant weights $\lambda_\tau$,
Hermitian matrices $G_\tau$ with nonzero determinant and full
$R$-lattices $M_\tau\subset K^n$. Write $V=VV(K,\lambda)$,
$P_G=\mathrm{pairK}(\lambda,G)$ and
$M_\tau^\#=\{x:\bar y^tG_\tau x\in R\text{ for all }y\in M_\tau\}$.
Use $\nabla(M),\Delta(M)$ for the two transported lattices restricted
to $V$. The two conclusions in \path{NabDelDual}
(\path{CoefModLemmas.lean}:143--148) are equalities of $R$-submodules of $V$:
\[
 \nabla(M)^{\vee,P_G}=\Delta(M^\#),\qquad
 \Delta(M)^{\vee,P_G}=\nabla(M^\#).
\]
The dual on the left always tests the \emph{first} argument against
the specified lattice. We prove both equalities at each finite place and
then intersect.

For a sesquilinear $h$ and a finitely generated $R$-module $T$ in its
first argument, suppose $x\in V$ and the module $D$ is characterized
on $V$ by $x\in D\iff h(T,x)\subset R$. For each finite prime $w$,
\[
 x\in D_w\quad\Longleftrightarrow\quad h(T_{\bar w},x)\subset R_w .
\]
If $s x\in D$ for $s\notin w$, and $t y\in T$ for $t\notin\bar w$,
then $h(ty,sx)=\bar t s h(y,x)\in R$, with $\bar t s\notin w$.
Conversely the local pairing condition holds on generators $y_j$ of
$T$. Choose $s_j\notin w$ clearing $h(y_j,x)$ and put $s=\prod_j s_j$.
Then $h(T,sx)\subset R$ by conjugate linearity and $sx\in V$,
so $sx\in D$. The empty product also works. For every $R$-submodule
$A$ inside a $K$-space one has $A=\bigcap_w A_w$: if $x\notin A$,
the proper ideal $\{s:sx\in A\}$ is contained in a maximal ideal
$w$, at which $x\notin A_w$. This applies also to $R\subset K$.
Hence $h(T,x)\subset R$ is equivalent to all the tests
$h(T_{\bar w},x)\subset R_w$; this last equivalence follows simply
by localizing each scalar and also holds without finite generation.
See \path{mem_of_forall_locP}, \path{pair_denom},
\path{mem_localization_dual} and \path{pair_integral_iff}
(\path{TransportedWeylLatticeHermitianDuality.lean}:29--44,70--99,101--123).

Fix $w$. Each $M_\tau^\#$ is full and finitely generated:
$y\mapsto\bar y^t G_\tau$ takes generators of $M_\tau$ to a
finite spanning family of row vectors over $K$. Its $R$-span is
therefore between $dR^n$ and $d^{-1}R^n$ for some nonzero rational
integer $d$ (clear the coordinates and a spanning inverse).
Taking integral dot-product duals gives the same sandwich for
$M_\tau^\#$, which is finitely generated since $R$ is Noetherian.
This supplies the lattice clause of \path{dualFr_lattice_and_hdual}
(\path{DualLatticeIsLatticeAndHermitianDual.lean}:92--147,173--175).
Choose $B_\tau$ with columns an $R_{\bar w}$-basis of
$(M_\tau)_{\bar w}$ and put $B'_\tau=(B_\tau^*G_\tau)^{-1}$.
Since $B_\tau^*G_\tau B'_\tau=1$, the tests against the columns of
$B_\tau$ read off the coordinates $(B'_\tau)^{-1}x$. By the
localized dual test these coordinates belong to $R_w$ precisely when
$x\in (M_\tau^\#)_w$. Thus $B'_\tau$ is a basis there, including
at primes fixed by the involution. This proves the frame step in
\path{mem_colLat_dual}, \path{isBasisAt_hdual} and \path{exists_dual_frame}
(\path{TransportedWeylLatticeHermitianDuality.lean}:162--201,204--210).

Take in turn $(T,T')=(\nabla,\Delta)$ and $(\Delta,\nabla)$ in the
standard polynomial model. Write $A=T(M)$ and $A'=T'(M^\#)$ for
the transported polynomial modules. The standard $T$ is finitely
generated by the integral bases above
(\path{std_fg_span}, \path{StdWeylLatticesFgAndSpanCoeffSpace.lean}:108--121).
Every element of $A$ belongs to $V$: the preceding localization
argument gives this for $\nabla(M)$, and
$\Delta_{\rm std}\subset\nabla_{\rm std}$ since its divided
words are integral polynomials in the same coefficient space.
Transport preserves inclusions by the defining local tests
(\path{transLat_mono}, \path{delStd_le_nabStd},
\path{CoefModLemmas.lean}:248--257).
Thus restricting to $V$ loses no test vectors.

For $x\in V$ set $x_0=\mathrm{act}(B'^{-1})x\in V$.
The localization formulas just established give
$A_{\bar w}=\mathrm{act}(B)T_{\bar w}$ and
$A'_w=\mathrm{act}(B')T'_w$. The covariance identity is
\[
 P_G(\mathrm{act}(B)y,x)
   =P_{B^*GB'}(y,x_0)=P_1(y,x_0).
\]
Indeed $x=\mathrm{act}(B')x_0$; adjointness of row substitution for
the factorial monomial pairing moves the first block substitution
to the second, and the Gram determinant twist transforms by
$\overline{\det\mathrm{Tw}(B)}\det\mathrm{Tw}(B')$, precisely
the two scalar factors in the actions. This is \path{pairTransport}
and its restriction \path{pairK_apply}
(\path{CoefModLemmas.lean}:477--504).
Therefore $P_G(A_{\bar w},x)\subset R_w$ is equivalent to
$P_1(T_{\bar w},x_0)\subset R_w$, hence to $x_0\in T'_w$ by
the localized dual test applied to standard Weyl duality, and
finally to $x\in A'_w$. Intersecting over all $w$ gives
$P_G(A,x)\subset R\iff x\in A'$ in $V$. The two choices of
$(T,T')$ prove both equalities. See \path{dual_of_local} and
\path{GlobalTheory.TransDual.nabDelDual}
(\path{TransportedWeylLatticeHermitianDuality.lean}:218--274,278--309).

For the Z4.4 numerator, \path{GlobalTheory.PeriodClassSum.pairBDualK}
sets the star on $K$ to the identity and $G_\tau=1$. The vector
dual then becomes the ordinary dot-product dual. Symmetry of the
bilinear normalized pairing identifies $P_1(y,x)$ with
$\mathrm{pairBK}(x,y)$; the two transported equalities give
the two integral tests. \path{PeriodClassSum.pairBDual} applies
this to the scalar extensions of the class lattice
(\path{IsClassSumPairingFormsLat.lean}:181--197).
\path{PairAtClass.pair_integral} uses the first test on the integral
class component and the supplied opposite row to get its coefficient
in $R$ (\path{ClassPairingIntegralFormRowIsIntegral.lean}:49--57).
This is the route from the registered standard Weyl and transported
localization theorems into the Z4.4 root. The public alias
\path{GlobalTheory.T.transported_weylLattice_hermitianDuality}
(\path{TransportedWeylLatticeHermitianDualityStatement.lean}:14--15)
states the same two equalities.

